\chapter{Anti-GBM Disease}
\index{Anti-GBM Disease}
\label{sec:Anti-GBM Disease (Goodpasture Syndrome)}

\section{Overview}
Anti-glomerular basement membrane disease, also known as Goodpasture syndrome, is a rare, life-threatening small-vessel vasculitis characterized by autoantibodies directed against the non-collagenous domain of the $\alpha$3 chain of type IV collagen ($\alpha$3(IV)NC1), found in the basement membranes of glomeruli and lung alveoli, with an incidence of 0.5--1 per million and a bimodal age distribution (young adults 20--30 years with male predominance, and older adults 60--70 years with female predominance). The condition involves anti-GBM antibodies binding to the target antigen in the GBM and alveolar basement membrane, activating complement and recruiting neutrophils, leading to crescentic glomerulonephritis and alveolar bleeding, with genetic susceptibility strongly associated with HLA-DRB1*1501.
\section{Causes}
\section{Risk Factors}

\begin{itemize}[noitemsep]
  \item Genetic predisposition and family history
  \item Advancing age
  \item Lifestyle factors (diet, exercise, smoking, alcohol)
  \item Environmental exposures
  \item Underlying medical conditions
  \item Medication use
\end{itemize}
\section{Signs and Symptoms}

\subsection{Common symptoms}
\begin{itemize}[noitemsep]
  \item You may experience rapidly progressive glomerulonephritis (hematuria, proteinuria, oliguria, rising creatinine, crescentic GN on biopsy), with or without lung involvement (hemoptysis, shortness of breath, alveolar bleeding---seen in 60--80\%, more common in smokers and young patients). 
  \end{itemize}

\subsection{Emergency warning signs}
\begin{itemize}[noitemsep]
  \item Severe or worsening symptoms of anti-gbm disease require immediate medical evaluation
\end{itemize}
\section{Progression and Stages}
The condition may progress through different stages of severity over time. Early stages often involve mild or subtle symptoms that may be overlooked. Moderate stages involve more noticeable symptoms affecting daily function. Advanced stages may involve significant impairment and complications.
\section{Complications}
\begin{itemize}[noitemsep]
  \item Disease progression leading to worsening symptoms
  \item Secondary infections or injuries
  \item Side effects of treatment
  \item Reduced quality of life
  \item Emotional and psychological impact
\end{itemize}
\section{Diagnosis}
Doctors diagnose this based on detection of circulating anti-GBM antibodies by ELISA and kidney biopsy showing linear IgG deposition along the GBM on immunofluorescence (the pathognomonic finding). Comprehensive medical history and physical examination Laboratory tests to assess organ function and rule out other conditions Imaging studies to visualize affected structures Specialized testing depending on the affected system Consultation with relevant specialists Monitoring of disease progression over time

\begin{itemize}[noitemsep]
  \item Comprehensive medical history and physical examination
  \item Laboratory tests to assess organ function
  \item Imaging studies to visualize affected structures
  \item Specialized testing depending on the affected system
  \item Consultation with relevant specialists
\end{itemize}
\section{Prevention}
\begin{itemize}[noitemsep]
  \item Maintain a healthy lifestyle with balanced diet and exercise
  \item Avoid known risk factors
  \item Attend regular medical check-ups for early detection
  \item Follow recommended screening guidelines
\end{itemize}
\section{Treatment Options}
Treatment typically involves urgent: high-dose corticosteroids combined with cyclophosphamide and plasmapheresis (daily for 7--14 days to remove pathogenic antibodies). Medications to manage symptoms and address underlying mechanisms Lifestyle modifications and self-care measures Physical therapy and rehabilitation Surgical intervention when indicated Regular monitoring and follow-up care Patient education and support services

\begin{itemize}[noitemsep]
  \item Medications to manage symptoms and address underlying mechanisms
  \item Lifestyle modifications and self-care measures
  \item Physical therapy and rehabilitation
  \item Surgical intervention when indicated
  \item Regular monitoring and follow-up care
\end{itemize}
\section{Living with It}
\begin{itemize}[noitemsep]
  \item Follow your treatment plan as prescribed
  \item Maintain regular follow-up with your healthcare provider
  \item Make necessary lifestyle adjustments
  \item Seek support from family, friends, or support groups
  \item Stay informed about your condition
\end{itemize}
\section{Outlook}
\begin{itemize}[noitemsep]
  \item Your outlook depends on the degree of kidney impairment at presentation
  \item If dialysis-dependent at diagnosis, kidney recovery is rare.
\end{itemize}
